\scanpage{250}
Voici un résultat \deleted{plus} profond, fondamental \deleted{\unsure{dans le cas|dans le}} en géométrie analytique :

\uline{Lemme} 1.3. (Grauert). Soit $f:X\unsure{\text{---}|\longrightarrow}Y$ une application propre d'espaces analytiques sur $\mathbb C$, $F$ un faisceau coh. sur $X$. Alors

(i) les $R^if(\unsure{F|B})$ sont cohérents

(ii) Soit $y\in Y$, \deleted{Alors le système}\inserted{$Y_n=\operatorname{Specan}(\mathcal O_y/\mathfrak m_y^{n+1})$}, $X_n=X\times_Y Y_n$, $F_n=F\otimes_Y Y_n$, alors le système projectif $(H^i(X_n,F_n))_n$ satisfait M.L. et l'hom canonique
\[
{R^if(F)_y}^{\wedge}\longrightarrow\varprojlim_n H^i(X_n,F_n)
\]
est un isom.

\uline{Corollaire} 1.4. Supposons que $y$ soit un point isolé du supp $R^if(F)$, \deleted{Alors} i.e que $R^if(F)_y$ soit artinien. Alors on a un isom
\[
R^if(F)_y\longrightarrow\varprojlim_n H^i(X_n,F_n)
\]
\mspage{249}
